\subsection{The relative correction is \texorpdfstring{$O(k^{-1})$}{O(1/k)} in the class with a gap}\label{sec:B3T06}
Retain the notation of Section~\ref{sec:B3T05} and set $q=1-p$.
On this small branch, Jacobi's formula from Section~\ref{sec:B3T01} takes the form
\[
 (\nabla\ell)_e=\frac{6p_e}{q_e(1-4p_e)}
 -\frac12\left[
 \frac{q_e^2+3p_e^2}{(1-4p_e)^2}A_eGA_e^T
 -\sum_{d\ni e}G_{dd}\right].
\]
The final sum includes only the two diagonals and excludes the RC constraints.
The expression also follows by differentiating the local factor and $S_{\rm gf}$. Indeed, $\partial_{p_e}\tau_e=(q_e^2+3p_e^2)/(1-4p_e)^2$, while the RC diagonal block of $S$ remains identically zero.
Section~\ref{sec:B3T05} gives $A_eGA_e^T=4+O(k^{-1})$ and $\sum_{d\ni e}G_{dd}=2+O(k^{-1})$.
Moreover, $p_e=O(k^{-1})$, so
\[
 \boxed{\nabla_p\ell=-\mathbf1+O_\infty(k^{-1}).}
\]
The constant term lies in an RC normal direction and contributes no effective correction to the probabilities.

\paragraph{Relative response.}
Using gauge-invariant forms, the Schur complement of the block $N$ gives
\[
 R=T-TAG A^TT.
\]
Since $\|P^{-1}T\|_\infty\le2$, $\|A\|_\infty=4$, and $\|A^TT\|_\infty=\max_\ell\sum_{e\in\ell}\tau_e\le2$ for large $k$,
\[
 \|P^{-1}R\|_\infty
 \le2(1+8\|G\|_\infty)=O_{\delta_0}(1).
\]
The constant vector $\mathbf1$ belongs to the normal space of the RC constraints, so $R\mathbf1=0$. Section~\ref{sec:B3T01} and the preceding two estimates give
\[
 \boxed{\delta=P^{-1}R(\nabla\ell+\mathbf1)=O_\infty(k^{-1}).}
\]
In particular, $1+\delta_e>0$ for sufficiently large $k$ in this class. The proof requires no additional lower bound on $p_e$. It does require $p_e>0$ and the uniform vacancy and weighted gap bounds of Section~\ref{sec:B3T05}.

\paragraph{First-order cancellation in the derivative term.}
The exact identity $\sum_{c\in r}p_c\delta_c=0$ from Section~\ref{sec:B3T01} implies
\[
 \sum_{c\in r}p_c\log(1+\delta_c)
 =-\frac12\sum_{c\in r}p_c\delta_c^2+O(k^{-3})=O(k^{-2}).
\]
Thus the first-order contribution of the derivative term cancels. For the actual insertion defect we must still add $r_c$, giving
\[
 \mathbb E_p\epsilon
 =\mathbb E_p r-\tfrac12\mathbb E_p\delta^2+O(k^{-3}).
\]
Section~\ref{sec:macro} controls $r_c$. The result holds uniformly in the specified class. We make no claim that every legal residual board in the queens problem belongs to this class.
